Put $X=L-3>0$, $Y=R-3>0$. The conditions are
\[
Y+6=aX,\qquad X+6=bY
\]
for positive integers $a,b$. They cannot both equal $1$.

If $a=1$, then $X=Y+6$ and $(b-1)Y=12$, so $Y$ is a positive divisor of $12$. The resulting totals $X+Y+6=2Y+12$ are $14,16,18,20,24,36$. The case $b=1$ gives the same totals by symmetry.

If $a,b\ge2$, the inequalities $2X\le Y+6$, $2Y\le X+6$ imply $X,Y\le6$. Testing these small positive values under the two divisibility conditions yields
\[
(X,Y)=(1,1),(2,2),(2,4),(3,3),(4,2),(6,6),
\]
giving totals $8,10,12,18$. The distinct possible totals are therefore $8,10,12,14,16,18,20,24,36$, whose sum is $\boxed{158}$.
